% Generated draft: requires mathematical and source review.
\documentclass[11pt]{article}
\usepackage[margin=25mm]{geometry}
\usepackage{fontspec}
\setmainfont{Latin Modern Roman}
\usepackage{amsmath,amssymb,mathtools,unicode-math}
\setmathfont{Latin Modern Math}
\usepackage{enumitem,needspace,xurl,hyperref}
\setlength{\parindent}{0pt}
\setlength{\parskip}{5pt}
\setlength{\emergencystretch}{3em}
\newcommand{\N}{\mathbb{N}}
\newcommand{\Z}{\mathbb{Z}}
\newcommand{\Q}{\mathbb{Q}}
\newcommand{\R}{\mathbb{R}}
\newcommand{\C}{\mathbb{C}}
\newcommand{\F}{\mathbb{F}}
\newcommand{\A}{\mathbb{A}}
\newcommand{\PP}{\mathbb{P}}
\newcommand{\E}{\mathbb{E}}
\newcommand{\eps}{\varepsilon}
\newcommand{\dd}{\,\mathrm{d}}
\begin{document}
\section*{UnsolvedMath ID: 800764}
\textbf{Source problem code:} catalogue entry
\subsection{Definitions and mathematical statement}
Let \(\mathcal L_{\mathrm{ring}}=\{0,1,+,-,\cdot\}\) be the first-order language of rings, and let \(K\) be an infinite field. A subset of \(K\) is cofinite if its complement in \(K\) is finite. For an integer \(m\geq 0\), an \(\mathcal L_{\mathrm{ring}}\)-formula \(\varphi(x;y_1,\ldots,y_m)\), and a tuple \(\mathbf a\in K^m\), put
\[
A_{\varphi,\mathbf a}=\{b\in K:K\models\varphi(b;\mathbf a)\}.
\]
When \(m=0\), this means a formula with no parameters.

The supplied source calls \(K\) minimal if every definable subset of \(K\) is finite or cofinite, but does not explicitly specify whether parameters from \(K\) are allowed. To make this qualification explicit, for \(I\in\bigl\{\{0\},\Z_{\geq 0}\bigr\}\), let \(M_I(K)\) denote the condition
\[
\text{for every }m\in I,\ \varphi(x;y_1,\ldots,y_m),\ \mathbf a\in K^m,
\quad A_{\varphi,\mathbf a}\text{ is finite or cofinite in }K.
\]
Here \(I=\{0\}\) expresses parameter-free definability, whereas \(I=\Z_{\geq0}\) allows arbitrary finite tuples of parameters from \(K\).

The supplied conjecture, attributed to Podewski, is that every minimal field is algebraically closed. For the intended parameter convention \(I\), its precise formulation is: for every infinite field \(K\),
\[
M_I(K)\ \Longrightarrow\ 
\left[
\begin{gathered}
\text{for every integer }n\geq1\text{ and all }c_0,\ldots,c_n\in K
\text{ with }c_n\neq0,\\
\text{there exists }b\in K\text{ such that }
\sum_{j=0}^{n}c_jb^j=0.
\end{gathered}
\right]
\]
The supplied record also states that algebraically closed fields are minimal. Which parameter convention the source intends remains to be confirmed; the two displayed choices are alternative readings, not an assertion that they are equivalent.
\subsection{Short English statement}
Podewski's conjecture says that an infinite field in which every definable subset is finite or has finite complement must be algebraically closed. The supplied wording leaves the use of parameters in definitions unspecified, so the draft makes both possible conventions explicit. The record labels the conjecture open; that status has not been independently verified.
\subsection{Sources}
\begin{itemize}
\item Local UnsolvedMath registry, permanent ID \(800764\), exact source problem code \texttt{catalogue entry}; supplied complete source record, upstream revision \texttt{d371ee917ac9647fbb06cd6fc529deec6f3dda82}.
\item Bary-Soroker and Fehm, \emph{Open Problems in the Theory of Ample Fields} (2011), Conjecture 7, \url{https://arxiv.org/abs/1106.1310}. The supplied record lists Bary-Soroker and Fehm as proposers and attributes the conjecture to Podewski.
\item Supplied provenance: accessed 2026-10-16; extraction method \texttt{arxiv\_latex}. This is source-provided access information, not a claim of an independent visit or verification. The supplied registry status is \emph{open}.
\end{itemize}
\end{document}
